\documentclass[10pt,a4paper]{amsart}
\usepackage[margin=19mm]{geometry}
\usepackage{amsmath,amssymb,mathtools}
\usepackage[scr=rsfs,cal=euler]{mathalfa}
\usepackage{xfrac}
\usepackage[unicode,hidelinks,bookmarks=false]{hyperref}
\newcommand{\ie}{i.e.\ }
\newcommand{\cf}{cf.\ }
\newcommand{\resp}{respectively\ }
\newcommand{\Subsection}{\subsection}
\providecommand{\nobreakdash}{\nobreak\hbox{-}}
\setlength{\emergencystretch}{2em}
\multlinegap0pt
\usepackage{amssymb}
\usepackage{xcolor}
\usepackage{csquotes}
\newtheorem{lemma}{Lemma}
\newcommand{\ov}[1]{\overline{#1}}
\DeclareMathOperator{\rk}{rk}
\DeclareMathOperator{\tr}{Tr}
\title{Chern classes of Laughlin bundles\\ on the quasihole moduli space}
\author{Florent Dupont (IRMA), Semyon Klevtsov (IRMA)}
\date{Source: arXiv:2605.18089v2 (CC BY 4.0)}
\begin{document}
\maketitle

\noindent\emph{Source and licence.} Chern classes of Laughlin bundles\\ on the quasihole moduli space. Version arXiv:2605.18089v2 (CC BY 4.0), distributed by arXiv under the Creative Commons Attribution 4.0 International licence, http://creativecommons.org/licenses/by/4.0/, as recorded in the arXiv licence metadata for this exact version and shown on the abstract page https://arxiv.org/abs/2605.18089v2. Full licence notice: \texttt{licenses/arxiv-2605.18089-CC-BY-4.0.txt}.\par\smallskip

\noindent\textit{Editorial scope.} Counted unit: the article's lemma lemma (source0.tex lines 919--921). The article's complete original proof of it is delivered with it (source0.tex lines 923--956, one \texttt{proof} environment of 34 lines). No mathematics is written, repaired or re-derived: the statement and the proof are the article's own lines, in the article's own notation, and no retained line is substituted or re-flowed.  No further source-local context is needed: every cross-reference of every retained line resolves inside the two delivered spans.  Nothing was omitted from any retained span, so the package carries no omission record.  No retained line contains a citation, so the wrapper carries no bibliography and no bibliography entry.  No retained line contains a figure environment, an image-inclusion command or drawing code, so no image is delivered and the package declares no asset dependency.  Editorial formatting only, in addition: an AMS \texttt{amsart} preamble that restates the article's own declarations and macros used by the retained lines, namely the environments \texttt{lemma} on separate counters, together with the standard packages those lines need.  The retained environments print in this document as Lemma 1 in order of appearance, whereas the article prints the same items with the numbering of its own full document.  The complete native source of the article as distributed in the arXiv e-print for this exact version is bundled unchanged as \texttt{sources/200834/source0.tex}.
\begin{lemma}
  Let $H'=Hh$. In cohomology $[\tr e^{-\frac{i}{2\pi}\partial_m \ov{\partial}_m \log H'}]=\rk(V)$.
\end{lemma}

\begin{proof}
The first term in the Taylor expansion of the exponential gives $\rk(V)$ when taking the trace, and we have to show that $\forall k>0$, $[\tr \left( (\partial_m \ov{\partial}_m \log H')^k\right)]=0$.

Set $$\chi=\tr \left( \left(\partial_m \ov{\partial}_m \log H'\right)^{k-1} \ov{\partial}_m \log H' \right).$$ 
From the previous Lemma it follows that
under $w_\gamma \to w_\gamma+1$ we have $H' \mapsto PH'P^{-1}$ and under $w_\gamma \to w_\gamma+\tau$, $H' \mapsto QH'Q^{-1}$, where $P$ and $Q$ are the following constant matrices
\[
P=
\begin{pmatrix}
1 & 0 & 0 & \cdots & 0 \\
0 & q & 0 & \cdots & 0 \\
0 & 0 & q^2 & \cdots & 0 \\
\vdots & \vdots & \vdots & \ddots & \vdots \\
0 & 0 & 0 & \cdots & q^{b-1}
\end{pmatrix},\quad\quad\quad Q= 
\begin{pmatrix}
0 & 0 & \cdots & 0 & 1 \\
1 & 0 & \cdots & 0 & 0 \\
0 & 1 & \cdots & 0 & 0 \\
\vdots & \vdots & \ddots & \vdots & \vdots \\
0 & 0 & \cdots & 1 & 0
\end{pmatrix},
\]
where recall $q=e^{\frac{2\pi i}{b}}$.
Upon conjugation of $H'$ by any invertible constant matrix $U$, $\chi$ is invariant. Indeed,
\begin{align*}
    &\tr \left( \left(\partial_m \ov{\partial}_m (\log (UH'U^{-1})\right))^{k-1} \ov{\partial}_m (\log (UH'U^{-1}))\right)\\
    &=\tr \left( \left(\partial_m \ov{\partial}_m (U\log H'U^{-1})\right)^{k-1} \ov{\partial}_m (U\log H'U^{-1})\right)\\
    &=\tr \left(  \left(U(\partial_m \ov{\partial}_m \log H' )U^{-1}\right)^{k-1} U\ov{\partial}_m (\log H')U^{-1}\right)\\
    &=\tr \left( U \left(\partial_m \ov{\partial}_m \log H' \right)^{k-1} \ov{\partial}_m \log H' U^{-1}\right)\\
    &=\chi
\end{align*} and thus $\chi$ is a smooth form defined on $S^m C$. Furthermore, $d\chi=\tr \left(\partial_m \ov{\partial}_m \log H'\right)^k$ hence the result.

\end{proof}

\end{document}
